
\section*{2. Polynômes de Snapper}

\subsection*{2.1.}
Rappelons d'abord la définition et quelques propriétés élémentaires des applications polynômes :

\textbf{Définition 2.1.1.} \uline{Soit} $f$ \uline{une application d'un groupe abélien} $A$ \uline{dans un groupe abélien} $B$\uline{. On dit que} $f$ \uline{est une} application polynôme\uline{, s'il existe, pour toute suite} $x_1,\ldots,x_r$ ($r\geq1$\uline{) d'éléments de} $A$\uline{, une famille} $(y(x_1,\ldots,x_r;k_1,\ldots,k_r))_{k_1,\ldots,k_r\geq0}$ \uline{d'éléments de} $B$\uline{, presque tous nuls, et satisfaisant à la condition suivante : Pour toute suite} $n_1,\ldots,n_r$ \uline{d'entiers, on a}
\[
f(n_1x_1+\cdots+n_rx_r)=
\sum_{k_1,\ldots,k_r\geq0}(n_1^{k_1}\cdots n_r^{k_r})
y(x_1,\ldots,x_r;k_1,\ldots,k_r).
\]
\uline{On dit que} $f$ \uline{est de degré} $\leq m$\uline{, si quels que soient} $r\in\mathbb Z$ \uline{et} $x_1,\ldots,x_r\in A$\uline{, on a}
\[
y(x_1,\ldots,x_r;k_1,\ldots,k_r)=0
\]
\uline{pour} $k_1+\cdots+k_r>m$.

\textbf{Lemme 2.1.2.} \uline{Tout homomorphisme de groupes abéliens est une application polynôme de degré} $\leq1$\uline{. Si} $h=f\circ g$\uline{, où} $f$ \uline{est une application polynôme de degré} $p$ \uline{et} $g$ \uline{une application polynôme de degré} $q$\uline{, alors} $h$ \uline{est une application polynôme de degré} $\leq pq$.

Démonstration triviale.

\textbf{Lemme 2.1.3.} \uline{Soit}
\[
0\longrightarrow G'\longrightarrow G\xrightarrow{j}G''\longrightarrow0
\]
\uline{une suite exacte de groupes abéliens telle que} $G'$ \uline{est divisible par}
